\section{Experiments}
\label{sec:testing}
 
This section describes the benchmark design and evaluation methodology used to assess Hypothetica's originality assessment system across a diverse set of research idea inputs. The evaluation is designed to test the full pipeline under realistic retrieval conditions rather than only measuring isolated model behavior. Each benchmark case is evaluated against one evidence source, passed through the same retrieval, reranking, parsing, and scoring stages used by the application, and then mapped to one of three originality labels.

The methodology has three goals. First, it ensures that the benchmark is balanced across corpora, originality labels, and input lengths, so that aggregate accuracy cannot be explained by class imbalance or by unusually verbose inputs. Second, it aligns each idea with the writing style of its target corpus, because scientific papers, patents, and software repositories expose different kinds of evidence. Third, it defines a strict labeling scheme before evaluation, which makes the results interpretable as a test of originality assessment rather than a loose relevance-matching exercise.
 
\subsection{Benchmark Dataset Overview}
 
We constructed three independent benchmark datasets, one per retrieval corpus, totaling 135 labeled cases (45 per corpus). Each dataset was designed to reflect the natural idea style of its corresponding corpus, and all cases were manually authored and labeled. Rather than sampling from existing labeled repositories or relying on automated generation, we deliberately crafted each idea to probe specific behaviors of the originality assessment pipeline -- including its ability to distinguish genuine conceptual innovation from standard domain extensions.
 
Each corpus contains exactly equal representation across the three classes (15 cases per class per corpus), ensuring that overall accuracy is not dominated by any single class and that per-class metrics are individually meaningful. A dataset skewed toward \textit{already\_exists} -- the natural outcome of random sampling from any real corpus -- would inflate overall accuracy without providing meaningful signal about novelty detection, which is the primary use case of the system.
 
\subsection{Idea Length Stratification}
 
A key design dimension in this benchmark is \textbf{idea length}. Within each class and each corpus, ideas are evenly distributed across three length tiers: \textit{short} (one to two sentences), \textit{medium} (three to four sentences), and \textit{long} (five or more sentences), with five cases per tier per class. This stratification tests whether the pipeline's performance is sensitive to the amount of detail provided in the input description.
 
The motivation for this design is that the pipeline processes the enriched idea description at multiple stages -- query variant generation, candidate selection, and per-criterion LLM scoring all depend on the textual content of the input. A short idea provides less signal for retrieval and a narrower context for the scoring agents. A long idea gives the retrieval layer more angles to search from and gives the Layer~1 agent more material to compare against each candidate. By fixing class labels across length tiers, we can isolate the effect of description length from the effect of originality class, and measure how much the pipeline benefits from additional input detail.
 
In real use, the follow-up question agent and idea enrichment step partially compensate for short inputs by eliciting additional context before retrieval begins. The benchmark deliberately tests raw idea text without this enrichment step, which means the length effect measured here represents a lower bound on performance for short ideas in production.
 
\subsection{Corpus-Specific Idea Design}
 
A core design principle was stylistic alignment: ideas in each dataset were written to match the natural framing conventions of the corpus against which they would be retrieved.
 
\textbf{OpenAlex (Scientific Literature).} Ideas were framed as research proposals and methodological contributions, consistent with the academic writing style of OpenAlex-indexed papers. Descriptions emphasize problem formulation, proposed mechanisms, and expected contributions -- the vocabulary a researcher would use when describing a novel method or system to peers.
 
\textbf{GitHub (Software Repositories).} Ideas were framed as practical software projects, consistent with the repository descriptions and README summaries found on GitHub. Descriptions are implementation-oriented, referencing specific technologies, frameworks, and functional behaviors rather than abstract mechanisms.
 
\textbf{Patents (Google Patents).} Ideas were framed as system-level inventions, consistent with the technical and functional language of patent abstracts. Descriptions emphasize novel system architectures, operational paradigms, and distinguishing technical claims rather than experimental contributions or code-level implementations.
 
This stylistic alignment was intentional: a retrieval system indexing academic papers will surface different candidates for the same underlying idea than one indexing patents or repositories. By matching framing to corpus, we ensured that retrieval quality -- and therefore scoring -- reflects realistic deployment conditions rather than a vocabulary mismatch artifact.
 
\subsection{Labeling Scheme: The Novelty Ladder}
 
All cases were labeled under a strict three-tier scheme, which we refer to as the novelty ladder:
 
\textbf{Already Exists.} The idea is widely known, commonly implemented, or straightforwardly retrievable in the target corpus. A case receives this label if a competent practitioner in the field would immediately recognize it as existing work. Representative examples include:
 
\begin{itemize}
    \item \textit{[GitHub, Testing, short]} ``A JavaScript testing setup using Jest and React Testing Library for unit tests, component tests, mocks, and coverage reports.'' (score: 14)
    \item \textit{[OpenAlex, AI/ML, short]} ``A deep residual network trained on ImageNet that uses skip connections between convolutional layers to enable stable training of very deep image classification models.'' (score: 9)
    \item \textit{[Patents, Healthcare, long]} ``An electronic health record system that stores patient data and allows secure sharing between healthcare providers. The system maintains medical history, diagnoses, medications, lab results, imaging reports, and clinical notes in a centralized or interoperable digital format. Access controls and audit logs are used to protect sensitive patient information.'' (score: 17)
\end{itemize}
 
\textbf{Incremental.} The idea shares its core concept with existing work but introduces a clear, limited improvement -- a better component, an added feature, a new application domain, or a modest architectural modification. The key criterion is that the improvement is a natural extension of prior work rather than a new mechanism. Combinations of the form ``existing method X applied to domain Y'' are classified as incremental regardless of whether the specific combination has been published. Representative examples include:
 
\begin{itemize}
    \item \textit{[GitHub, Blockchain, long]} ``A smart contract auditing tool that combines symbolic execution results with LLM-generated explanations of potential vulnerabilities. The symbolic engine identifies suspicious execution paths while the LLM converts findings into developer-friendly explanations and suggested fixes.'' (score: 41)
    \item \textit{[OpenAlex, Biomedical, long]} ``A multimodal survival prediction model for lung adenocarcinoma that combines whole-slide histopathology images with bulk RNA sequencing profiles. A cross-attention module learns relationships between visual morphology and molecular signatures.'' (score: 38)
    \item \textit{[Patents, Manufacturing, long]} ``A predictive maintenance system that learns machine-specific degradation curves by fine-tuning a cross-factory foundation model on local sensor history. The base model captures general failure patterns across many factories, while local adaptation accounts for equipment-specific wear behavior.'' (score: 56)
\end{itemize}
 
\textbf{Novel.} The idea introduces a genuinely new mechanism, control paradigm, or representational approach that cannot be reduced to a recombination of existing components. Novel cases must change how a problem is fundamentally approached, not merely what is applied to it. We explicitly excluded superficial novelty -- ideas that appear original through buzzword recombination (e.g., ``ML + X,'' ``blockchain-based Y'') but do not introduce a new underlying principle. Representative examples include:
 
\begin{itemize}
    \item \textit{[GitHub, Testing, long]} ``A software testing framework that generates tests by simulating adversarial future maintainers who misunderstand the codebase in different ways. The framework models likely human misreadings of function names, abstractions, and architectural boundaries, then creates tests that protect against bugs introduced by those misunderstandings.'' (score: 82)
    \item \textit{[OpenAlex, AI/ML, long]} ``A continual learning system where forgetting is treated as a structured compression process. After each task, the model reorganizes old knowledge into fractal attractors in weight space, leaving sparse regions available for new tasks rather than preventing forgetting through regularization or replay.'' (score: 83)
    \item \textit{[Patents, Smart Home, long]} ``A home automation system that constructs a causal model of resident behavior from passive sensor streams and uses counterfactual reasoning to infer unstated goals. It asks what outcome the resident was likely trying to achieve under different device states, then proposes new automation rules that would have produced that outcome with less effort.'' (score: 93)
\end{itemize}
 
This distinction between incremental and novel was the most consequential design decision. A strict interpretation of novelty -- requiring a new mechanism or paradigm rather than a new combination -- produces a labeling scheme that approximates how a domain expert or peer reviewer would assess originality, rather than how a keyword-matching system would.
 
\subsection{Deliberate Inclusion of Borderline Cases}
 
Each dataset includes a proportion of deliberately borderline cases -- ideas that sit close to the incremental/novel boundary and would reasonably invite disagreement among human annotators. This design choice reflects the reality of originality assessment: the boundary between a highly creative extension and a genuinely new mechanism is often ambiguous in practice, and a useful assessment tool must behave meaningfully in this region rather than only on clear-cut cases. These cases also serve as a natural stress test for the system's sensitivity near decision boundaries, and their presence partly explains the lower per-class accuracy on the incremental label observed across all three corpora.
 
\subsection{Default Pipeline Parameters}
 
Table~\ref{tab:params} lists the fixed parameter values used in all experiments. Retrieval and selection parameters control how many candidates reach the scoring phases; scoring parameters determine how per-candidate similarity scores are combined into the final originality verdict.
 
\begin{table}[width=\textwidth,cols=3,pos=ht]
    \caption{Default pipeline parameters used during benchmark evaluation.}
    \label{tab:params}
    \begin{tabular*}{\textwidth}{@{\extracolsep{\fill}}p{3.4cm}p{4.5cm}p{5.8cm}@{}}
        \toprule
        Parameter & Value & Role \\
        \midrule
        Query variants & 4 & Diversified retrieval formulations per idea \\
        Papers per query variant & 150 & Initial retrieval depth per variant \\
        Retained hits per variant & 40 & Candidate retention before deduplication \\
        Embedding top-$k$ & 100 & First-stage semantic retrieval cutoff \\
        Rerank top-$k$ & 20 & Cross-encoder reranking cutoff \\
        Final analyzed candidates & 5 & Documents forwarded to Layer~1 scoring \\
        Likert mapping & $\{1,2,3,4,5\}\rightarrow\{0,0.25,0.5,0.75,1.0\}$ & Criterion score normalization \\
        Criterion weights $(w_p,w_m,w_d,w_c)$ & $(0.15,0.30,0.10,0.45)$ & Problem, method, domain, and contribution weighting \\
        Global aggregation $(w_{\max},w_{\text{mean}})$ & $(0.70,0.30)$ & Max-mean similarity aggregation across candidates \\
        Originality curve exponent $p$ & 1.5 & Nonlinear conversion from similarity to originality \\
        Sentence overlap thresholds & 0.70 / 0.40 & High- and medium-overlap annotation cutoffs \\
        Default score thresholds & $<$40 / 40--69 / $\geq$70 & Already\_exists, incremental, and novel labels before corpus calibration \\
        \bottomrule
    \end{tabular*}
\end{table}
 
 
